\documentclass[12pt]{article}
\usepackage[utf8]{inputenc}
\usepackage[T1]{fontenc}
\usepackage{hyperref}
\title{Upper Ta'oih Language Resources for K-12 Curriculum Development}
\author{Compiled by Claude Sonnet 5 (claude-sonnet-5) for Andr\'as Kornai}
\date{}
\begin{document}
\maketitle
\section{Introduction}
Upper Ta'oih (ISO 639-3 tth) is a Katuic (Austroasiatic) variety of the Ta'Oi dialect continuum, spoken in Salavan and Sekong provinces of Laos and in Thua Thien-Hue province of Vietnam. Ta'Oi is not usually written; the Vietnamese and Lao varieties are only partly mutually intelligible, and the group has several ISO codes (tth, tto, irr, oog), so sources must be checked carefully.
\section{Primary Educational Resources}
No primers or school materials specific to Upper Ta'oih were found in search results.
\section{Linguistic and Reference Materials}
\begin{itemize}
\item Wikipedia, Ta'Oi language (covers the whole group): \url{https://en.wikipedia.org/wiki/Ta%27Oi_language}
\item Omniglot page on Ta'Oi: \url{https://mail.omniglot.com/writing/taoi.htm}
\end{itemize}
\section{Religious and Community Resources}
\begin{itemize}
\item Scripture Earth listing for tth: \url{https://www.scriptureearth.org/00eng.php?iso=tth}
\item Joshua Project, Upper Ta Oi in Laos: \url{https://joshuaproject.net/index.php/people_groups/15260/LA}
\item Joshua Project, Ta'oih in Vietnam: \url{https://joshuaproject.net/index.php/people_groups/15260/VM}
\end{itemize}
The contents of the Scripture Earth page were not opened and are unverified.
\section{Digital and Audio Resources}
None found in search results.
\section{Curriculum Development Context}
Upper Ta'oih is not known to be taught in schools. Literacy materials would need to be built from scratch or from SIL work on related Katuic varieties, and the choice of script (Lao or Vietnamese-based Latin) would depend on the country.
\end{document}